\paragraph{从已知比例走到状态方程}
若已知质量比例，可把任意共同总质量 $m$ 暂作计数尺度：$m_i=w_im$、$n_i=w_im/\mu_i$。总摩尔数为
\[
 n=m\sum_i\frac{w_i}{\mu_i},\qquad
 \overline\mu=\frac mn=\left(\sum_i\frac{w_i}{\mu_i}\right)^{-1}.
\]
共同的 $m$ 消去，说明结果只依赖比例。再求 $x_i=(w_i/\mu_i)/\sum_j(w_j/\mu_j)$，最后才可用 $p_i=x_ip$。这一顺序同时核对了摩尔质量单位与比例类型。
道尔顿分压关系在这里以理想、同温、各组分共享同一体积的混合物为前提。若分子间相互作用不能忽略，或有凝结、化学反应改变组分粒子数，不能把初始配比不加判断地当作当前气相摩尔分数。本文的 $\mu_i$ 表示摩尔质量，其含义由本节定义限定。
